\documentclass[11pt,letterpaper]{article}
\usepackage[T1]{fontenc}
\usepackage{amsmath,amsthm}
\usepackage{newtxtext,newtxmath}
\usepackage[margin=1in,headheight=15pt]{geometry}
\usepackage{microtype,mathtools,booktabs,array,longtable}
\usepackage[dvipsnames]{xcolor}
\usepackage{enumitem,fancyhdr,xurl}
\usepackage[most]{tcolorbox}
\usepackage{hyperref}
\hypersetup{colorlinks=true,linkcolor=MidnightBlue,citecolor=MidnightBlue,
 urlcolor=MidnightBlue,pdftitle={Finite Separators and Discriminant Factorizations for Sextic Block Resolvents},
 pdfauthor={Research article prepared for Vladimir Reshetnikov}}
\usepackage[nameinlink,noabbrev]{cleveref}
\setlength{\parskip}{4pt}
\setlength{\parindent}{1.1em}
\setlist{nosep,leftmargin=*}
\pagestyle{fancy}
\fancyhf{}
\fancyhead[L]{\small\scshape Finite sextic separators}
\fancyhead[R]{\small ProveIt research}
\fancyfoot[C]{\thepage}
\renewcommand{\headrulewidth}{0.3pt}
\numberwithin{equation}{section}
\setcounter{tocdepth}{1}
\newtheorem{theorem}{Theorem}[section]
\newtheorem{lemma}[theorem]{Lemma}
\newtheorem{proposition}[theorem]{Proposition}
\newtheorem{corollary}[theorem]{Corollary}
\theoremstyle{definition}
\newtheorem{definition}[theorem]{Definition}
\newtheorem{example}[theorem]{Example}
\newtheorem{question}[theorem]{Question}
\theoremstyle{remark}
\newtheorem{remark}[theorem]{Remark}
\DeclareMathOperator{\disc}{disc}
\DeclareMathOperator{\Gal}{Gal}
\DeclareMathOperator{\lc}{lc}
\DeclareMathOperator{\sgn}{sgn}
\newcommand{\QQ}{\mathbb Q}
\newcommand{\ZZ}{\mathbb Z}
\newcommand{\CC}{\mathbb C}
\newcommand{\NN}{\mathbb N}
\newcommand{\Ptwo}{\mathcal P_2}
\newcommand{\Pthree}{\mathcal P_3}
\newcommand{\Edges}{\mathcal E}
\newcommand{\Rtwo}{R_{2,f}}
\newcommand{\Rthree}{R_{3,f}}
\newcommand{\repo}{\texttt{ProveIt}}
\newcommand{\src}[1]{\texttt{\nolinkurl{#1}}}
\newcommand{\snapshot}{\texttt{e2b1f016a94102663f12b970e6dd434229f90d01}}
\newtcolorbox{resultbox}{enhanced,breakable,colback=MidnightBlue!3,
 colframe=MidnightBlue!50,boxrule=0.5pt,arc=1pt,left=9pt,right=9pt,top=7pt,bottom=7pt}

\begin{document}
\begin{titlepage}
\vspace*{0.35in}
{\color{MidnightBlue}\large\scshape Algebra \; / \; Invariant theory \; / \; Verified computation\par}
\vspace{0.25in}
{\Huge\bfseries Finite Separators and\par Discriminant Factorizations\par for Sextic Block Resolvents\par}
\vspace{0.25in}
{\Large A sharp 196-point guarantee and a bounded replacement\par for ProveIt's separating-invariant search\par}
\vspace{0.35in}
{\large Prepared for Vladimir Reshetnikov\par}
{\large September 28, 2026\par}
\vspace{0.28in}
\begin{abstract}
The sextic radical-solvability development in \repo\ uses terminating but
unbounded search to specialize two block-partition descriptors. We replace
that search, at the level of an ordinary mathematical proof, by a single
explicit parameter range. For every separable sextic in characteristic zero,
one of the integers $1,\ldots,196$ separates all fifteen pair partitions and
all ten triple partitions simultaneously. The underlying descriptors compress
to quadratics and linears. Their collisions are controlled by three invariant
polynomials $Q_f,U_f,L_f$, with degrees at most $30,120,45$, and we prove
\[
 \disc_Z\Rtwo=\disc(f)^6Q_f^6U_f^2,
 \qquad
 \disc_Z\Rthree=\disc(f)^3L_f^2.
\]
For the generic sextic, the three factors are irreducible and have the stated
exact degrees. Consequently 196 is optimal for a guarantee valid for every
set of distinct nonzero complex test parameters, although optimality of the
initial integer interval is not established. A split integer example gives a
lower bound of eight for that latter problem. We derive a coefficient-only
primitive-recursive decision architecture, a height-only separator requiring
no search, a precise repository integration plan, and further research
questions. Exact symbolic checks and a modular B\'ezout certificate accompany
the proofs.
\end{abstract}
\vfill
\begin{resultbox}
\textbf{Status and scope.} This is an unrefereed research article with complete
written proofs and reproducible exact checks. The new results have not been
formalized or run in Lean. General decidability, including polynomial-time
decidability of solvability by radicals, is classical; it is not claimed as a
new discovery. The specific separator and collision-factorization results are
proposed contributions relative to the inspected repository. Historical
priority is not certified by the literature check.
\end{resultbox}
\end{titlepage}
\tableofcontents
\clearpage

\section{The repository question and the contribution}
\label{sec:scope}

The starting point is a concrete implementation boundary, not an assertion
that sextic solvability is an unsolved problem. In the pinned \repo\ snapshot,
\src{SexticRadicalDecision.lean} proves a total recursive criterion for integer
sextics. Its introductory comment explicitly does not claim primitive
recursiveness, because the separating-invariant search uses unbounded
minimization. The supporting file \src{SexticIrreducibleDecision.lean}
constructs pair and triple parameters using \src{Nat.find}. The file
\src{SexticSeparatingInvariants.lean} proves that suitable parameters exist by
representing each block by its root polynomial and each partition by the
polynomial of those block polynomials \cite{repo,repo-decision,repo-search,repo-invariants}.

That is a particularly useful research target: the existing semantic
architecture already identifies the right invariant and its correctness
condition. What is missing from this route is a small explicit bound and a
sufficiently transparent explanation of the exceptional specializations.
Our answer is stronger than merely bounding an enumeration of pairs of
natural numbers. A one-dimensional curve works for every distinct root
tuple, its failure set has three explicit components, and its generic number
of failures can be determined exactly.

\subsection{What is classical, inherited, and proved here}

The equivalence between radical solvability and solvability of the Galois
group is classical. So is the sextic criterion using the two imprimitive
wreath-product groups of degrees fifteen and ten in the corresponding
coset actions; see, for example, Boswell and Glasser \cite{boswell-glasser}.
Landau and Miller proved polynomial-time decidability of radical solvability
in 1985 \cite{landau-miller}. Thus even a primitive-recursive decision for
sextics is not, as an abstract existence theorem, new.

The partition descriptors and the strategy of separating their specialized
values are inherited from the inspected \repo\ source. More broadly, the
construction and use of relative invariants are central to computational
Galois theory; Fieker and Kl\"uners discuss general methods and the practical
importance of structured invariants \cite{fieker-klueners}.

The claims established in this article are the explicit compression,
the special curve $(u,v)=(t,2t^2)$, the $196$-point uniform bound, the
three-factor description of its obstruction, the exact discriminant
identities, and generic irreducibility and sharpness. Their proofs are given
below rather than inferred from computations. The bounded decision
architecture is an application to the repository's mathematical gap, not a
replacement claim for the stronger classical complexity theorem.

The literature check consulted the primary sources just cited and the
pinned source files. It does not establish that no equivalent identity has
appeared elsewhere under different coordinates or resolvent conventions.
Accordingly, ``contribution'' here means a result developed and proved for
this report, with a demonstrated use relative to this repository; it does
not mean a certified first-in-literature result.

\subsection{Main results at a glance}

The root tuple is always assumed distinct. The polynomials $h_P$ and $k_T$
are defined explicitly in \cref{sec:compression}; their resolvents have degrees
fifteen and ten in $Z$.

\begin{resultbox}
\textbf{Uniform separator.} Among any 196 distinct nonzero elements of a
characteristic-zero field containing the roots, there is a parameter $t$
for which both maps $P\mapsto h_P(t)$ and $T\mapsto k_T(t)$ are injective.
In particular, testing $1,\ldots,196$ always succeeds.

\textbf{Exact obstruction.} There is a nonzero polynomial
$E_f=Q_fU_fL_f$ of degree at most 195, with coefficients polynomially
expressible in the sextic coefficients, such that $E_f(t)\ne0$ is exactly
simultaneous separation. Its three factors have generic degrees $30,120,45$
and are generically irreducible.

\textbf{Sharpness and implementation.} The number 196 is optimal for
arbitrary complex test sets. For the initial integer interval the proven
bounds are $8\le N_{\mathrm{int}}\le196$. The explicit finite range removes
unbounded minimization from a coefficient-only sextic decision architecture.
\end{resultbox}

\section{Notation and complete partition descriptors}
\label{sec:descriptors}

Let $K$ be a field of characteristic zero, let
$\alpha_1,\ldots,\alpha_6\in K$ be distinct, and put
\begin{equation}
 f(X)=\prod_{i=1}^6(X-\alpha_i)
     =X^6-e_1X^5+e_2X^4-e_3X^3+e_4X^2-e_5X+e_6.
 \label{eq:sextic}
\end{equation}
If the coefficients originally lie in a smaller field, take $K$ to contain
a splitting field. Our arguments are algebraic and do not depend on an
ordering, numerical approximation, or choice of complex branches.

Write $[6]=\{1,\ldots,6\}$ and
$\Edges=\binom{[6]}2$. Let $\Ptwo$ be the set of partitions of $[6]$ into
three unordered two-element blocks, and let $\Pthree$ be the set of
partitions into two unordered three-element blocks. Their cardinalities are
\begin{equation}
 |\Ptwo|=\frac{6!}{(2!)^3 3!}=15,
 \qquad |\Pthree|=\frac{\binom63}{2}=10.
\end{equation}
We call members of $\Ptwo$ \emph{matchings}. When two matchings have no common
block, we call them \emph{block-disjoint}; they still use the same six roots.

For a block $B\subseteq[6]$, define its root polynomial and the complete
partition descriptor by
\begin{equation}
 p_B(u)=\prod_{i\in B}(u-\alpha_i),
 \qquad D_{\mathcal B}(u,v)=\prod_{B\in\mathcal B}(v-p_B(u)).
 \label{eq:descriptor}
\end{equation}
These are the ordinary polynomial counterparts of the descriptors in
\src{SexticSeparatingInvariants.lean} \cite{repo-invariants}. The distinction
between the variables $u$ and $v$ matters: $u$ enters the root polynomials,
whereas $v$ records the unordered collection of blocks.

\begin{lemma}[A descriptor determines the partition]
\label{lem:complete}
For distinct roots, the map $\mathcal B\mapsto D_{\mathcal B}(u,v)$ is
injective on either partition family.
\end{lemma}
\begin{proof}
First, the map $B\mapsto p_B(u)$ is injective. Indeed, the roots of $p_B$
among $\alpha_1,\ldots,\alpha_6$ are precisely the roots indexed by $B$.
Next regard $D_{\mathcal B}$ as a monic polynomial in $v$ over the field
$K(u)$. Its distinct linear factors are $v-p_B(u)$. Equality of two such
products implies equality of their sets of roots in $K(u)$, so it implies
equality of the sets of block polynomials. Their injectivity recovers the
blocks and hence the partition. This proof also explains why repeated
sextic roots must be excluded at this stage.
\end{proof}

The descriptors therefore solve the information-loss problem before
specialization. The remaining task is to evaluate them at a pair $(u,v)$
without causing two distinct descriptors to become equal. Arbitrary fixed
evaluations need not preserve injectivity. We next expose the unusually
small amount of partition-dependent data in the sextic case.

\section{Compression to quadratics and linears}
\label{sec:compression}

\subsection{Pair partitions}

For a pair $b=\{i,j\}$, write
$s_b=\alpha_i+\alpha_j$ and $p_b=\alpha_i\alpha_j$. For a matching $P$, set
\begin{align}
 A_P&=\sum_{b\in P}p_b,\notag\\
 B_P&=\sum_{\substack{b,c\in P\\b\ne c}}s_b p_c,\notag\\
 C_P&=\sum_{\{b,c\}\in\binom P2}p_b p_c,
 &h_P(t)&=A_Pt^2+B_Pt-C_P.
 \label{eq:h}
\end{align}
The sum defining $B_P$ is over ordered pairs of different blocks, so it has
six summands before expanding the $s_b$.

\begin{proposition}[Pair descriptor expansion]
\label{prop:pair-expansion}
For every matching $P$,
\begin{align}
 D_P(u,v)={}&v^3-(3u^2-e_1u+A_P)v^2\notag\\
 &+\bigl(3u^4-2e_1u^3+(e_2+A_P)u^2-B_Pu+C_P\bigr)v-f(u).
 \label{eq:pair-descriptor}
\end{align}
Consequently, if $K_2(t)=2t^6+2e_2t^4-f(t)$, then
\begin{equation}
 D_P(t,2t^2)=K_2(t)-2t^2h_P(t).
 \label{eq:pair-curve}
\end{equation}
The map $P\mapsto h_P(t)$ is injective as a map into $K[t]$.
\end{proposition}
\begin{proof}
Each block polynomial is $u^2-s_bu+p_b$. Their sum is
$3u^2-e_1u+A_P$. The sum of their pairwise products has leading coefficients
$3,-2e_1$, and its coefficient of $u^2$ is
$2A_P+\sum_{b<c}s_bs_c$. Since the global $e_2$ is
$A_P+\sum_{b<c}s_bs_c$, that coefficient is $e_2+A_P$.
The remaining two coefficients are $-B_P$ and $C_P$. The product of all
three block polynomials is $f(u)$. Expanding the monic cubic in $v$ gives
\cref{eq:pair-descriptor}; substituting $u=t,v=2t^2$ gives
\cref{eq:pair-curve}.

If two $h$ polynomials agree, their coefficients give equality of $A,B,C$.
The displayed descriptor formula then gives equal descriptors. Apply
\cref{lem:complete}.
\end{proof}

\subsection{Triple partitions}

For a three-element block $B$, let $e_j(B)$ denote the elementary symmetric
polynomial of degree $j$ in its three roots. For $T\in\Pthree$, define
\begin{equation}
 a_T=\sum_{B\in T}e_2(B),\qquad
 b_T=\sum_{B\in T}e_3(B),\qquad
 k_T(t)=a_Tt-b_T.
 \label{eq:k}
\end{equation}

\begin{proposition}[Triple descriptor expansion]
\label{prop:triple-expansion}
For every triple partition $T$,
\begin{equation}
 D_T(u,v)=v^2-(2u^3-e_1u^2+a_Tu-b_T)v+f(u).
 \label{eq:triple-descriptor}
\end{equation}
On $(u,v)=(t,2t^2)$ this has the form
\begin{equation}
 D_T(t,2t^2)=K_3(t)-2t^2k_T(t),
 \quad K_3(t)=4t^4-4t^5+2e_1t^4+f(t).
 \label{eq:triple-curve}
\end{equation}
The map $T\mapsto k_T(t)$ into $K[t]$ is injective.
\end{proposition}
\begin{proof}
The sum of the two monic cubic block polynomials is
$2u^3-e_1u^2+a_Tu-b_T$, and their product is $f(u)$. Expanding their
quadratic descriptor proves the formula. Equality of $k$ polynomials gives
equality of the two descriptor-dependent coefficients, and then
\cref{lem:complete} proves injectivity.
\end{proof}

In particular every difference $h_P-h_Q$ for $P\ne Q$ is a nonzero
polynomial of degree at most two; every difference $k_T-k_{T'}$ is nonzero
of degree at most one. A first root-counting bound would therefore allow
$2\binom{15}{2}+\binom{10}{2}=255$ bad parameters. The next two sections
remove repeated collision factors and reduce this to 195.

\begin{remark}[The zero parameter and the source coordinate order]
The compressed families make sense and may separate at $t=0$, but the
original descriptors on the curve all collapse within each family there,
because the multiplier $2t^2$ vanishes. We use positive integer tests when
transferring results to the repository. In its \src{Fin 2} convention,
\src{X 0} is the outer descriptor variable $v$ and \src{X 1} is $u$.
Thus the required parameter is $x(0)=2t^2$, $x(1)=t$, not the reverse
\cite{repo-invariants}.
\end{remark}

\section{The two types of pair-partition collision}
\label{sec:pair-collisions}

Two different matchings can share at most one block, since two common pairs
force the remaining pair to agree. There are exactly 45 unordered pairs of
matchings sharing a block: choose the common pair in 15 ways, then choose
two of the three matchings on the four remaining indices. The other
$\binom{15}{2}-45=60$ pairs are block-disjoint. Their union is a single
alternating six-cycle, since a cycle of length two would be a common block.

For a root pair $e=\{i,j\}$, define
\begin{equation}
 q_e(t)=t^2+s_et-p_e
       =t^2+(\alpha_i+\alpha_j)t-\alpha_i\alpha_j.
 \label{eq:q}
\end{equation}

\begin{lemma}[Common-block collision factor]
\label{lem:shared}
If distinct matchings $P,Q$ share the pair $e$, then
\begin{equation}
 h_P(t)-h_Q(t)=(A_P-A_Q)q_e(t),\qquad A_P-A_Q\ne0.
 \label{eq:shared}
\end{equation}
\end{lemma}
\begin{proof}
Let the four remaining roots be $a,b,c,d$. Their two-block pairing affects
$A$ through one of
\[
 z_1=ab+cd,\qquad z_2=ac+bd,\qquad z_3=ad+bc.
\]
The products of the two remaining pair products are always $abcd$.
The contribution to $B$ involving only those two blocks is always the
third elementary symmetric function of $a,b,c,d$. Therefore changing the
pairing changes $B$ by $s_e$ times the change in $A$, and changes $C$ by
$p_e$ times the change in $A$. Substitution in \cref{eq:h} proves the
factorization.

The three differences are
\begin{equation}
 z_1-z_2=(a-d)(b-c),\quad
 z_1-z_3=(a-c)(b-d),\quad
 z_2-z_3=(a-b)(c-d).
 \label{eq:four-root-differences}
\end{equation}
Distinctness of the roots makes every one nonzero.
\end{proof}

Thus the 45 shared-block comparisons do not contribute 90 unrelated
exceptional parameters: they contribute roots of only fifteen quadratics.
The 60 block-disjoint comparisons each still contribute a quadratic.
There are at most $2(15+60)=150$ pair-family failures. Some of these
quadratics can have lower degree or share roots after specializing the
sextic, so this is a uniform upper bound rather than a claim about every
input.

\section{Triple collisions as exchanges of two roots}
\label{sec:triple-collisions}

For disjoint root pairs $e,e'$, set
\begin{equation}
 \ell_{e,e'}(t)=(s_e-s_{e'})t-(p_e-p_{e'}).
 \label{eq:ell}
\end{equation}
Interchanging $e,e'$ changes only the sign. There are
$\binom64\cdot3=45$ unordered choices of two disjoint pairs.

\begin{lemma}[Exchange identity]
\label{lem:exchange}
For the triple partitions
\[
T=\{\{a,b,c\},\{d,e,f\}\},\qquad
T'=\{\{a,b,d\},\{c,e,f\}\},
\]
where letters now denote root values,
\begin{equation}
 k_T(t)-k_{T'}(t)
   =(c-d)\bigl((a+b-e-f)t-(ab-ef)\bigr).
 \label{eq:exchange}
\end{equation}
Every unordered pair of distinct triple partitions occurs exactly once in
this manner, up to the irrelevant orientations, from an unordered pair of
disjoint root pairs $\{\{a,b\},\{e,f\}\}$.
\end{lemma}
\begin{proof}
Subtracting the internal pair sums gives
$a_T-a_{T'}=(c-d)(a+b-e-f)$, while subtracting the two triple products gives
$b_T-b_{T'}=(c-d)(ab-ef)$. This proves the formula.

For the combinatorial assertion, two distinct partitions can be oriented so
that one block from each has two elements in common. Call these $a,b$;
call the exchanged elements $c,d$; the other two elements are $e,f$.
Conversely, two disjoint pairs and their two complementary elements specify
exactly the two partitions in the display. Simultaneously switching the
chosen blocks or reversing the exchange does not change the unordered data.
This establishes the bijection, with 45 elements on either side.
\end{proof}

\begin{lemma}[No identically vanishing exchange factor]
\label{lem:ell-nonzero}
For disjoint pairs of distinct roots, $\ell_{e,e'}$ is not the zero
polynomial.
\end{lemma}
\begin{proof}
If it vanished identically, the two pair sums and the two pair products
would agree. The two monic quadratics $X^2-s_eX+p_e$ and
$X^2-s_{e'}X+p_{e'}$ would agree, contradicting their disjoint root sets.
\end{proof}

The triple-family failure set is therefore the union of the roots of 45
nonzero linears, and has cardinality at most 45. A nonzero constant factor
simply contributes no root.

\section{Three invariant obstruction polynomials}
\label{sec:obstruction}

Fix total orders on the matchings and on the root pairs, solely to orient
differences in products. Define
\begin{align}
 Q_f(t)&=\prod_{e\in\Edges}q_e(t),\label{eq:Q}\\
 U_f(t)&=\prod_{\substack{P<Q\\P\cap Q=\varnothing}}(h_P(t)-h_Q(t)),
 \label{eq:U}\\
 L_f(t)&=\prod_{\substack{e<e'\\e\cap e'=\varnothing}}
             \ell_{e,e'}(t),\label{eq:L}\\
 E_f(t)&=Q_f(t)U_f(t)L_f(t).\label{eq:E}
\end{align}
Here $P\cap Q=\varnothing$ refers to their block collections. Reversing a
chosen orientation can multiply $U_f$ or $L_f$ by $-1$, which changes
neither the zero set nor any squared identity below.

It is not automatic that an oriented product of differences is symmetric;
the ordinary Vandermonde product is the familiar warning. The following
simple sign argument proves what is needed here.

\begin{lemma}[Graph-product sign test]
\label{lem:graph-sign}
Let a finite group act on a finite graph's vertices, permuting quantities
$w_v$. Form an oriented product $\prod_{\{v,w\}\in E}(w_v-w_w)$, with one
orientation for each edge. For a group element $\tau$ of order two, its
sign on this product is $(-1)^r$, where $r$ is the number of edges whose
endpoints are interchanged by $\tau$.
\end{lemma}
\begin{proof}
Edges exchanged with other edges occur in two-element orbits. The two
orientation signs in each such orbit multiply to $+1$, since applying
$\tau$ twice fixes each oriented difference. A fixed edge contributes $+1$
if its endpoints are individually fixed, and $-1$ if they are interchanged.
Multiplying these contributions gives the claim.
\end{proof}

\begin{proposition}[Symmetry, integrality, and exact failure test]
\label{prop:obstruction}
The expressions $Q_f,U_f,L_f$ are symmetric polynomials with integer
coefficients in the six roots and $t$. They are nonzero for every distinct
root tuple and satisfy
\begin{equation}
 \deg Q_f=30,\qquad \deg U_f\le120,\qquad \deg L_f\le45.
 \label{eq:degree-bounds}
\end{equation}
Their coefficients are universal integer polynomials in $e_1,\ldots,e_6$.
For every $t\in K$,
\begin{equation}
 E_f(t)\ne0
 \quad\Longleftrightarrow\quad
 \bigl(P\mapsto h_P(t)\bigr)\text{ and }
 \bigl(T\mapsto k_T(t)\bigr)\text{ are both injective}.
 \label{eq:exact-test}
\end{equation}
\end{proposition}
\begin{proof}
The product $Q_f$ is visibly symmetric. For $U_f$, use the graph with
matchings as vertices and block-disjoint pairs as edges. A transposition
of two root labels sends any matching either to itself or to a matching
sharing its untouched third block. Hence it interchanges the endpoints of
no edge in this graph. By \cref{lem:graph-sign}, it fixes $U_f$. Root
transpositions generate $S_6$, so $U_f$ is symmetric.

For $L_f$ use the graph on two-element subsets, joined when they are
disjoint, with vertex quantity $s_et-p_e$. A root transposition sends a
pair either to itself or to a pair sharing one root. It therefore reverses
no disjointness edge, so the same argument proves symmetry of $L_f$.
The fundamental theorem of symmetric polynomials over $\ZZ[t]$ now yields
the integer coefficient expressions in the $e_i$.

Every factor of $Q_f$ is monic quadratic. Every factor of $U_f$ is nonzero
by \cref{prop:pair-expansion}, and every factor of $L_f$ is nonzero by
\cref{lem:ell-nonzero}. This proves nonvanishing and the degree bounds.
Finally, \cref{lem:shared} accounts for all shared-block collisions,
\cref{eq:U} accounts for all block-disjoint ones, and
\cref{lem:exchange} accounts for all triple collisions. Their scalar
root-difference factors are nonzero. This proves \cref{eq:exact-test}.
\end{proof}

\begin{remark}[Homogeneity and scaling]
Assign weight one to every root and to $t$. Then $h_P,k_T,q_e,\ell_{e,e'}$
have respective total degrees $4,3,2,2$. Consequently $Q_f,U_f,L_f,E_f$
have total degrees $30,240,90,360$. Scaling all roots and $t$ by $c$
multiplies $E_f(t)$ by $c^{360}$. These are homogeneous degrees in the
roots and parameter together, not the parameter degrees in
\cref{eq:degree-bounds}.
\end{remark}

\section{Exact discriminant factorizations}
\label{sec:discriminants}

Define the compressed resolvents
\begin{equation}
 \Rtwo(Z,t)=\prod_{P\in\Ptwo}(Z-h_P(t)),\qquad
 \Rthree(Z,t)=\prod_{T\in\Pthree}(Z-k_T(t)).
 \label{eq:resolvents}
\end{equation}
Their coefficients are symmetric integer polynomials in the roots, so they
are coefficient-only constructions just as $E_f$ is. Let
\begin{equation}
 \Delta_f=\disc(f)=\prod_{1\le i<j\le6}(\alpha_i-\alpha_j)^2.
\end{equation}
For a monic polynomial with roots $z_1,\ldots,z_m$, we use the convention
$\disc=\prod_{i<j}(z_i-z_j)^2$.

\begin{theorem}[Discriminant identities]
\label{thm:discriminants}
As polynomial identities in the roots and $t$,
\begin{align}
 \disc_Z\Rtwo(Z,t)&=\Delta_f^6Q_f(t)^6U_f(t)^2,
 \label{eq:disc-pair}\\
 \disc_Z\Rthree(Z,t)&=\Delta_f^3L_f(t)^2.
 \label{eq:disc-triple}
\end{align}
In particular the formulas remain valid at nonseparating parameter values.
\end{theorem}
\begin{proof}
The pair-resolvent discriminant is the product of all 105 squared matching
differences. The 60 block-disjoint differences give $U_f^2$ directly.
Fix a common pair $e$. The three unordered comparisons between matchings
containing $e$ give, by \cref{lem:shared},
\[
 q_e(t)^6\prod_{1\le i<j\le3}(z_i-z_j)^2.
\]
By \cref{eq:four-root-differences}, the second factor is exactly the
discriminant of the monic quartic whose roots are the other four roots.
Multiplying over the fifteen choices of $e$ gives $Q_f^6$ times the
product of these fifteen quartic discriminants. A fixed root difference
$(\alpha_i-\alpha_j)^2$ occurs in precisely six of those quartics: choose
the excluded pair among the four other roots. Their product is therefore
$\Delta_f^6$, proving \cref{eq:disc-pair}.

For triples, \cref{lem:exchange} gives one squared linear factor for each
unordered pair of disjoint root pairs, producing $L_f^2$. The two leftover
roots are the exchanged roots. For a fixed exchanged pair there are three
pairings of the other four roots, so its squared difference appears three
times. The product of the root-difference factors is $\Delta_f^3$.
This proves \cref{eq:disc-triple}. No division was used, so these are
polynomial identities, not merely identities on the open set of good
parameters.
\end{proof}

\begin{corollary}[Discriminant square classes]
\label{cor:squareclass}
Over any characteristic-zero coefficient field, the pair-resolvent
discriminant is a square. If both $f$ and the triple resolvent are
separable, the triple-resolvent discriminant has the same square class
as $\Delta_f$.
\end{corollary}
\begin{proof}
The two formulas read
\[
 \disc_Z\Rtwo=(\Delta_f^3Q_f^3U_f)^2,
 \qquad
 \disc_Z\Rthree=\Delta_f(\Delta_f L_f)^2.
\]
In the second assertion all displayed factors are nonzero.
\end{proof}

There is a useful independent permutation explanation. A transposition of
root labels fixes three of the fifteen matchings and interchanges the
other twelve in six pairs; it acts evenly. It fixes four triple partitions
and interchanges the other six in three pairs; it acts oddly. Thus the
fifteen-point action has trivial sign character and the ten-point action
has the ordinary six-point sign character. The explicit factorizations
identify much more than this parity information: they also isolate every
collision factor and its multiplicity.

\section{Generic irreducibility of the three components}
\label{sec:irreducibility}

The degree bound is not merely the result of an inefficient factor count.
Generically all its factors are needed. In this section take
$x_1,\ldots,x_6,t$ to be algebraically independent over $\QQ$ and put
$S=\QQ[x_1,\ldots,x_6,t]$. Its symmetric subring is
\[
 S^{S_6}=\QQ[e_1,\ldots,e_6,t].
\]
The symbols $Q,U,L$ now denote the universal polynomials in this ring,
prior to any specialization of the sextic coefficients.

\subsection{Irreducibility of the individual root-coordinate factors}

\begin{lemma}[Primitive linear-polynomial test]
\label{lem:primitive}
Let $R$ be a unique factorization domain. If $A,B\in R$ are coprime and
$A\ne0$, then $Ay+B$ is irreducible in $R[y]$.
\end{lemma}
\begin{proof}
The polynomial is primitive, and is linear, hence irreducible, over the
fraction field of $R$. Gauss's lemma gives irreducibility over $R$.
\end{proof}

\begin{lemma}[The three kinds of factor are irreducible]
\label{lem:individual-irred}
In $S$, each $q_e$, each $\ell_{e,e'}$ for disjoint pairs, and each
$h_P-h_Q$ for block-disjoint matchings is irreducible.
\end{lemma}
\begin{proof}
For $q_e$, rename its two roots $a,b$ and write
\[
 q_e=a(t-b)+t(t+b).
\]
The coefficients as a linear polynomial in $a$ are coprime: substituting
$t=b$ in the second gives $2b^2\ne0$. Apply \cref{lem:primitive}.
For a linear exchange factor, rename its disjoint pairs $\{a,b\}$ and
$\{c,d\}$. Then
\[
 \ell=a(t-b)+(b-c-d)t+cd.
\]
Substituting $t=b$ in its constant coefficient gives
$(b-c)(b-d)\ne0$, so the same test applies. Adjoining unused independent
variables preserves these irreducibilities.

For the third type, all block-disjoint matching pairs are related by a
root permutation, because their union is an alternating six-cycle. It
suffices to consider
\[
 P=(ab)(cd)(ef),\qquad Q=(ac)(be)(df).
\]
Write $h_P-h_Q=aA+B$, where $A,B\in\QQ[b,c,d,e,f,t]$. This is linear in
$a$ because every monomial in \cref{eq:h} uses each root at most once.
It is homogeneous of total degree four; hence $A$ is homogeneous of
degree three and $B$ of degree four. At $t=0$, direct expansion gives
\begin{align}
 A_0&=-bcd+bce-bef+cdf,\notag\\
 B_0&=def(b-c).
 \label{eq:primitive-specialization}
\end{align}
The polynomial $A_0$ is divisible by none of $d,e,f,b-c$. Indeed its
respective substitutions at zero or at $b=c$ are
\[
 be(c-f),\quad cd(f-b),\quad bc(e-d),\quad c(e-d)(c-f),
\]
all nonzero. Therefore $\gcd(A_0,B_0)=1$.

Suppose an irreducible nonconstant polynomial $H$ divided both $A$ and
$B$. A factor of a homogeneous polynomial can be chosen homogeneous, so
$H$ is homogeneous. If $H|_{t=0}$ is nonzero, it has the same positive
total degree and divides both $A_0$ and $B_0$, a contradiction. If it is
zero, then $t$ divides $H$, so irreducibility makes $H$ associated to $t$;
but $A_0\ne0$. Thus $A,B$ are coprime. Apply
\cref{lem:primitive} with variable $a$ to finish the proof.
\end{proof}

The homogeneity in the last argument is essential: without it, coprime
specializations need not imply coprimality of the original polynomials.
Here it prevents a nonconstant common homogeneous factor from specializing
to a nonzero constant.

\subsection{Descent from root coordinates to coefficients}

\begin{lemma}[One-orbit product lemma]
\label{lem:orbit-product}
Let a finite group act on a polynomial ring over a field. Suppose a finite
collection of pairwise nonassociate irreducible polynomials forms a single
orbit up to units, and its product $F$ is invariant. Then $F$ is
irreducible in the invariant subring.
\end{lemma}
\begin{proof}
A factorization $F=AB$ in the invariant subring is also a factorization in
the ambient unique factorization domain. The irreducible factors of $A$
therefore form a subset of the factors of $F$. Invariance makes this subset
stable under the group. Transitivity forces it to be empty or the whole
set. In the first case $A$ is a unit, and in the second $B$ is a unit.
\end{proof}

\begin{theorem}[Generic three-factor theorem]
\label{thm:generic}
The universal polynomials $Q,U,L$ are pairwise nonassociate irreducible
polynomials in $\QQ[e_1,\ldots,e_6,t]$, and are irreducible in
$\QQ(e_1,\ldots,e_6)[t]$. Their degrees in $t$ are respectively
$30,120,45$. Their product $E$ is squarefree over the generic coefficient
field and has degree 195 and nonzero constant coefficient.
\end{theorem}
\begin{proof}
The individual factors are irreducible by \cref{lem:individual-irred}.
The symmetric group is transitive on root pairs, on unordered pairs of
disjoint root pairs, and on unordered pairs of block-disjoint matchings.
In each collection the individual factors are pairwise nonassociate.
For $q_e$, this follows from its monic leading term and its two root
indices. For $\ell_{e,e'}$, the coefficient of $t$ specifies the two
pairs with opposite signs, up to their interchange. For $h_P-h_Q$, the
coefficient of $t^2$ is $A_P-A_Q$: its six distinct root-product monomials,
with coefficients $+1$ or $-1$, specify the two block-disjoint matchings
up to interchange. Thus no further associations occur.

Each product is invariant by \cref{prop:obstruction}. The one-orbit
product lemma proves irreducibility in the symmetric subring. The leading
coefficient in $t$ of each individual $q$ is one, of each individual $h$
difference is a nonzero $A_P-A_Q$, and of each individual $\ell$ is a
nonzero difference of pair sums. The ambient ring is a domain, so the
three parameter degrees are exactly $30,120,45$. They are consequently
pairwise nonassociate. Their irreducibility and positive parameter degree
also imply that they are primitive over $\QQ[e_1,\ldots,e_6]$.
Gauss's lemma proves irreducibility over its fraction field.

In characteristic zero every irreducible polynomial of positive degree is
separable. The three are pairwise coprime there, so $E$ is squarefree and
has degree 195. Its constant coefficient is nonzero: every $q_e(0)$ and
$\ell_{e,e'}(0)$ is a nonzero polynomial, and every block-disjoint $h$
difference at zero is nonzero by the representative calculation
\cref{eq:primitive-specialization} and root permutation.
\end{proof}

This theorem is a statement about the generic coefficient field, not about
every numerical sextic. After specialization a factor can become
reducible, lose degree, or share roots with another factor. None of these
possibilities threatens the uniform separator theorem; they only reduce
the number of bad parameters. On the open coefficient locus
$\Delta_f\ne0$, the three polynomials describe the three irreducible
families of bad specializations over $\QQ$. The constructive factor proofs
also work after extending the constant field to any characteristic-zero
field.

\section{The finite separator theorem and its sharpness}
\label{sec:bound}

\begin{theorem}[Uniform simultaneous separation]
\label{thm:196}
For any six distinct roots in a characteristic-zero field, at most 195
parameters cause a collision in either compressed family. Therefore any
196 distinct nonzero test parameters contain a simultaneous separator.
The integers $1,\ldots,196$ always contain one, and this same parameter
separates both original descriptors on $(u,v)=(t,2t^2)$.
\end{theorem}
\begin{proof}
By \cref{prop:obstruction}, the bad parameters are exactly the roots of the
nonzero polynomial $E_f$, which has degree at most $30+120+45=195$.
A nonzero polynomial over a field has at most its degree many distinct
roots. Characteristic zero makes the positive integers $1,\ldots,196$
distinct and nonzero. Finally use \cref{eq:pair-curve,eq:triple-curve},
whose common multiplier is nonzero at these parameters.
\end{proof}

The same proof gives 151 test values for the pair family alone and 46 for
the triple family alone. A common parameter is preferable when a single
uniform interface is wanted; separate searches can stop sooner on some
inputs. There is no requirement that the least good pair parameter agree
with the least good triple parameter.

\begin{theorem}[Sharpness for arbitrary test sets]
\label{thm:sharp}
The constant 196 in the assertion about \emph{every} set of distinct
nonzero complex test parameters cannot be lowered. There are separable
sextics with exactly 195 distinct nonzero bad parameters. Such a sextic
can even be chosen split with integer roots.
\end{theorem}
\begin{proof}
Over the generic coefficient field, \cref{thm:generic} gives a degree-195
squarefree polynomial with nonzero constant coefficient. Thus its leading
coefficient, constant coefficient, and discriminant in $t$ are all nonzero.
Multiply these coefficient polynomials by $\Delta_f$ and substitute the
elementary symmetric functions of six independent root variables. The
result is still a nonzero polynomial in those variables, because the
symmetric-polynomial substitution is injective.

The integer lattice is Zariski dense in affine six-space: a nonzero
polynomial over $\QQ$ cannot vanish on all integer tuples, as follows by
induction on the number of variables and the one-variable root bound.
Choose an integer tuple where the displayed product is nonzero. Its roots
are distinct, and its $E_f$ has degree 195, no zero root, and no repeated
root. Over $\CC$ it has exactly 195 distinct nonzero roots. These are all
bad by \cref{eq:exact-test}; taking them as the test set disproves any
195-point guarantee of the stated kind.
\end{proof}

\subsection{An explicit sharpness certificate}

The existence proof above is entirely algebraic. The accompanying code also
gives a concrete witness with roots
\begin{equation}
 1,\ 4,\ 10,\ 23,\ 51,\ 109.
 \label{eq:sharp-roots}
\end{equation}
For the orientation used in the verifier, reduction modulo
$p=1{,}000{,}003$ gives
\begin{equation}
 \deg\overline E_f=195,\qquad
 \overline E_f(0)=521291,\qquad
 \lc(\overline E_f)=670728.
\end{equation}
The file \src{checks/sharpness_certificate.json} stores the coefficients of
$\overline E_f$ and polynomials $A,B$ satisfying the exact identity
\begin{equation}
 A\overline E_f+B\overline E_f'=1\quad\text{in }\mathbb F_p[t].
 \label{eq:bezout}
\end{equation}
The verifier checks primality of $p$ by trial division and checks this
identity by finite polynomial arithmetic. Degree preservation, the nonzero
constant coefficient, and \cref{eq:bezout} imply that the integer
polynomial $E_f$ has degree 195 and 195 distinct nonzero complex roots.
This is a reproducible computational certificate, not a Lean-checked one.
It is supplementary: \cref{thm:sharp} does not rely on its execution.

\subsection{The initial integer interval is a different problem}

Define $N_{\mathrm{int}}$ to be the least positive integer $N$ such that
$1,\ldots,N$ contains a simultaneous separator for every distinct complex
root tuple. The sharpness proof using complex roots of $E_f$ does
\emph{not} prove $N_{\mathrm{int}}=196$.

\begin{proposition}[An eight-point lower bound]
\label{prop:integer-lower}
For this initial-interval problem,
\begin{equation}
 8\le N_{\mathrm{int}}\le196.
\end{equation}
\end{proposition}
\begin{proof}
The upper bound is \cref{thm:196}. For the lower bound use the distinct
integer roots $(-3,0,1,2,3,4)$, whose polynomial is
\[
 X^6-7X^5+5X^4+55X^3-126X^2+72X.
\]
Direct substitution in \cref{eq:h,eq:k} gives the following exact numbers
of distinct values:
\begin{center}
\begin{tabular}{c*{8}{r}}
\toprule
$t$&1&2&3&4&5&6&7&8\\
\midrule
$|\{h_P(t)\}|$&12&13&13&14&15&15&15&15\\
$|\{k_T(t)\}|$&10&8&9&10&9&8&9&10\\
\bottomrule
\end{tabular}
\end{center}
Thus none of $1,\ldots,7$ works. For additional explicit checks, the sorted
pair values at $t=8$ are
\[
\begin{gathered}
-1128,-1044,-870,-792,-724,-486,-432,-424,\\
72,90,192,396,776,816,936,
\end{gathered}
\]
and the sorted triple values are
\[
 -32,-28,-22,-20,-14,1,32,46,92,160.
\]
They are all distinct within their respective families.
\end{proof}

\begin{proposition}[No universal single positive parameter]
\label{prop:no-single}
For each fixed positive integer $t_0$, there is a split separable integer
sextic for which $t_0$ fails to separate the pair family.
\end{proposition}
\begin{proof}
Take roots $0,t_0,2t_0,3t_0,4t_0,5t_0$. For the common pair
$\{2t_0,3t_0\}$, its quadratic satisfies
$q(t_0)=t_0^2+5t_0^2-6t_0^2=0$. Any two different matchings containing
that pair collide by \cref{lem:shared}.
\end{proof}
The example is reducible. It does not establish that a fixed parameter
must fail on some irreducible rational sextic; that is a separate question.

\section{Why the separated resolvents decide an irreducible sextic}
\label{sec:galois}

The group-theoretic criterion in this section is classical. We include its
proof to show exactly what the new separator supplies and what it does not.
It is the criterion underlying the degree-fifteen and degree-ten sextic
resolvents in the literature and the repository
\cite{boswell-glasser,repo-search}. Let $f\in\QQ[X]$ now be irreducible of
degree six, let $M$ be its splitting field, and view
$G=\Gal(M/\QQ)$ as a transitive subgroup of $S_6$.

\begin{lemma}[Solvable primitive groups have prime-power degree]
\label{lem:primitive-group}
A finite solvable group acting faithfully and primitively on a set with
more than one element has degree $p^r$ for some prime $p$.
\end{lemma}
\begin{proof}
Choose a minimal nontrivial normal subgroup $N$ of $G$. Because $G$ is
solvable, $N$ is elementary abelian. Here is the standard short justification:
a last nontrivial term of the derived series of $N$ is an abelian
characteristic subgroup of $N$, hence normal in $G$, so minimality makes
$N$ abelian. Its primary components are characteristic, so only one prime
$p$ occurs. The characteristic subgroup of elements killed by $p$ is
nontrivial and hence all of $N$.

The orbits of a normal subgroup form a $G$-invariant partition. Primitivity
and faithfulness force $N$ to be transitive: otherwise its orbits are
singletons, making $N$ act trivially. An abelian transitive permutation
group is regular. Indeed, an element fixing one point fixes every point
by commuting with all elements that carry it to another point. Faithfulness
then makes the stabilizer trivial. Thus the degree is $|N|=p^r$.
\end{proof}

\begin{proposition}[The two block systems are complete]
\label{prop:blocks-galois}
For a transitive subgroup $G\le S_6$, the following are equivalent:
$G$ is solvable; $G$ is imprimitive; $G$ preserves a partition in $\Ptwo$
or in $\Pthree$.
\end{proposition}
\begin{proof}
A solvable primitive group cannot have degree six by
\cref{lem:primitive-group}. Thus solvability implies imprimitivity.
For a transitive action, the blocks in a nontrivial block system have a
common size properly dividing six, so that size is two or three.
This proves the next implication.

The full stabilizer of a pair partition is $S_2\wr S_3$, an extension of
$S_2^3$ by $S_3$. The full stabilizer of a triple partition is
$S_3\wr S_2$, an extension of $S_3^2$ by $S_2$. Both are solvable,
since finite direct products and extensions of solvable groups are
solvable. Any subgroup preserving the corresponding partition is a
subgroup of its stabilizer and is therefore solvable.
\end{proof}

\begin{theorem}[Separated rational-root criterion]
\label{thm:criterion}
Let $f\in\QQ[X]$ be an irreducible monic sextic, and let $t\in\QQ$ be a
simultaneous separator. Then $f$ is solvable by radicals if and only if
at least one of $\Rtwo(Z,t)$ and $\Rthree(Z,t)$ has a rational root.
\end{theorem}
\begin{proof}
For a matching $P$ and $\sigma\in G$, coefficient equivariance gives
$\sigma(h_P(t))=h_{\sigma P}(t)$, since $t$ is rational. If $P$ is fixed
by $G$, then $h_P(t)$ is fixed by $G$ and lies in $\QQ$. Conversely,
if $h_P(t)$ is rational, then
$h_{\sigma P}(t)=h_P(t)$ for every $\sigma$, and separation gives
$\sigma P=P$. A rational root of the pair resolvent is one of its
explicit values, so it is equivalent to a $G$-fixed pair partition.
The same proof applies to $k_T$ and the triple resolvent.
Now use \cref{prop:blocks-galois} and the classical equivalence between
solvability by radicals and solvability of the Galois group.
\end{proof}

Separation is not a cosmetic condition in the converse. Without it,
$h_{\sigma P}(t)=h_P(t)$ would not imply $\sigma P=P$. A specialized value
could be fixed while its underlying partition moves among colliding
partitions. This is precisely the logical failure prevented by the
nonvanishing certificate $E_f(t)\ne0$.

For a monic integer sextic, both resolvents are monic integer polynomials
in $Z$. Any rational root is therefore an integer, so the final condition
can be decided using a finite integer interval. No splitting field or
numerical root approximation needs to occur in the returned-data algorithm.

\section{A bounded coefficient-only decision architecture}
\label{sec:algorithm}

\subsection{Computing the invariant coefficients without roots}

A root-coordinate definition can hide a computational oracle unless the
coefficient bridge is made explicit. Here the bridge is the constructive
fundamental theorem of symmetric polynomials. For each fixed power of $Z$
in \cref{eq:resolvents}, its coefficient is a fixed symmetric polynomial
in six roots with coefficients in $\ZZ[t]$. It can be reduced to a fixed
polynomial in $e_1,\ldots,e_6,t$ using integer arithmetic only.

For completeness, the reduction is finite in a particularly elementary
sense. Order root monomials lexicographically with
$x_1>\cdots>x_6$. The leading exponent vector of a symmetric polynomial
has the form $\lambda_1\ge\cdots\ge\lambda_6$. The product
\begin{equation}
 e_1^{\lambda_1-\lambda_2}e_2^{\lambda_2-\lambda_3}
 \cdots e_5^{\lambda_5-\lambda_6}e_6^{\lambda_6}
 \label{eq:symmetric-reduction}
\end{equation}
has exactly that leading root monomial, with coefficient one. Subtract its
multiple by the leading coefficient, which is a polynomial in $t$.
The leading monomial strictly decreases and the total root degree never
increases. There are only finitely many root monomials within the original
degree bound, so the process terminates with an expression over
$\ZZ[e_1,\ldots,e_6,t]$. This proof gives a finite symbolic construction,
not a choice of an unspecified algebraic root.

For the coefficient of $Z^{15-j}$ in $\Rtwo$, the root degree is at most
$4j$ and parameter degree at most $2j$. For the coefficient of $Z^{10-j}$
in $\Rthree$, they are at most $3j$ and $j$. Hence degrees at most 60 and
30 in the root variables suffice for the two complete coefficient tables.
These are fixed bounds, independent of the input height. The analogous
bound for $E$ is 360 in the root variables and 195 in $t$.

The article does not distribute expanded versions of these potentially
large universal tables. Their construction and their evaluation are
mathematically specified by the fixed finite products and the reduction
above. The supplied executable code is instead a verifier on exact root
tuples. This distinction is important when assessing what has and has not
been implemented.

\subsection{Monicization and bounded factor search}

Let the input be an integer polynomial
$g(X)=a_6X^6+\cdots+a_0$, with $a_6\ne0$. Define
\begin{equation}
 f(Y)=a_6^5g(Y/a_6)
     =Y^6+\sum_{i=0}^5 a_i a_6^{5-i}Y^i.
 \label{eq:monicization}
\end{equation}
This is monic integral, and its roots are $a_6$ times those of $g$.
Multiplication and division by the nonzero rational $a_6$ preserve the
radical closure of $\QQ$, so this transformation preserves all-root
radical solvability. The repository already provides this semantic bridge
\cite{repo-semantics}.

Write $H$ for the maximum absolute value of the six lower coefficients of
$f$, and let $R=1+H$. Every complex root of $f$ has absolute value at most
$R$. For example, if $|z|>1+H$, the geometric sum estimate gives
\[
 H\sum_{i=0}^5|z|^i<|z|^6,
\]
so the leading term cannot be canceled by the others.
If a monic degree-$d$ integer polynomial divides $f$, its roots are a
submultiset of these roots, and the coefficient at distance $j$ below its
leading term has absolute value at most
\begin{equation}
 \binom dj R^j.\label{eq:factor-bound}
\end{equation}
Gauss's lemma implies that rational reducibility of a monic integer
polynomial has a monic integer factor. Any proper factorization of a
sextic has a factor of degree at most three. Thus enumeration of the
finite coefficient boxes in \cref{eq:factor-bound}, for $d=1,2,3$,
completely decides reducibility.

A degree-two or degree-three factor leaves a quotient of degree at most
four; then every root is radical. A degree-one factor leaves a quintic,
whose radical solvability must actually be tested. In particular, one
must not accept every reducible sextic: a nonsolvable quintic multiplied
by a linear factor is the basic counterexample. The existing
primitive-recursive quintic decision can be reused
\cite{repo,repo-reducible}. Repeated factors are handled by these reducible
branches; the irreducible branch is automatically separable in
characteristic zero.

\subsection{The algorithm}

The following is a mathematical specification of a finite coefficient
algorithm. A default return on malformed input makes it a total Boolean
on seven encoded integers; the target predicate includes exact degree six.

\begin{resultbox}
\textbf{Bounded sextic radical criterion.}

\textbf{1.} Reject if $a_6=0$. Otherwise form the monic integer polynomial
$f$ in \cref{eq:monicization} and the root bound $R$.

\textbf{2.} Search the explicit boxes for a monic factor of degree two or
three. If one exists, accept.

\textbf{3.} Search the explicit interval for a linear factor. If one
exists, choose the first in a fixed order, form its quintic quotient,
and return the existing primitive-recursive quintic criterion.

\textbf{4.} The remaining input is irreducible. For $t=1,\ldots,196$,
evaluate the universal coefficient polynomials of the two compressed
resolvents. Choose the first $t$ for which both are squarefree, or,
equivalently, for which $E_f(t)\ne0$.

\textbf{5.} Accept exactly when one of the two resulting monic integer
resolvents has an integer root. For each resolvent, test all integers in
its Cauchy interval $[-(1+H_R),1+H_R]$, where $H_R$ is the maximum absolute
value of its lower coefficients.
\end{resultbox}

Squarefreeness in step 4 can be tested by a rational polynomial gcd with
the derivative, by a fixed-size resultant determinant, or by the
obstruction value itself. For an especially direct primitive-recursive
implementation, the fixed-size determinant avoids any concern about a
hidden unbounded gcd loop. The resolvent degrees are always fifteen and
ten.

\begin{theorem}[Bounded architecture and correctness]
\label{thm:primrec}
The coefficient procedure above has a primitive-recursive Boolean
realization and accepts exactly the integer inputs of degree six whose
complex roots are all expressible by radicals over $\QQ$.
\end{theorem}
\begin{proof}
Monicization is integer arithmetic with fixed exponents. The factor boxes
have primitive-recursive bounds, and polynomial division for their fixed
degrees is a finite coefficient computation. The correctness of their
branches follows from \cref{eq:factor-bound} and the degree arguments
above. The quintic subroutine is the established primitive-recursive
criterion, not an unbounded search for radical expressions.

On the irreducible branch, \cref{thm:196} guarantees that step 4 finds a
parameter. Each coefficient is the evaluation of a fixed integer
polynomial. Evaluation, fixed-size determinants, a loop of length 196,
and bounded selection are primitive recursive under the usual effective
encodings of integers and finite tuples. Step 5 is a bounded integer
search, and is complete because the resolvents are monic and the Cauchy
bound includes all their integer roots. Its correctness is
\cref{thm:criterion}. Composing these primitive-recursive operations proves
the claim.
\end{proof}

This theorem concerns a mathematical realization. It does not assert that
the existing Lean function has already been changed, that a \src{Primrec}
proof has been accepted by the kernel, or that the attached verifier is
this coefficient solver. A source-level integration must additionally
supply the coefficient-evaluation lemmas and bounded-selector code
discussed in \cref{sec:integration}.

Nor does the naive factor and integer-root enumeration above claim
polynomial bit complexity. Its bounds depend polynomially on coefficient
\emph{magnitudes}, which can still give exponential running time in their
bit lengths. The stronger general polynomial-time theorem is classical
\cite{landau-miller}; our objective is a small, transparent invariant
search with a tractable verification boundary.

\section{A separator determined by coefficient height alone}
\label{sec:height}

The finite search has only 196 candidates, but it is also possible to give
one coefficient-dependent parameter that is guaranteed to work without
any separation test. It is useful as a termination certificate or a
comparison construction, not as a recommendation for practical evaluation.

For a polynomial $P(t)=\sum c_it^i$ over $\CC$, write
$\|P\|_1=\sum|c_i|$. This norm is submultiplicative. Let $f\in\ZZ[X]$ be
monic and separable, put $R=1+H$ as above, and use $|\alpha_i|\le R$.
Then
\begin{align}
 \|q_e\|_1&\le1+2R+R^2\le4R^2,\notag\\
 \|h_P-h_Q\|_1&\le6R^2+24R^3+6R^4\le36R^4,\notag\\
 \|\ell_{e,e'}\|_1&\le4R+2R^2\le6R^2.
 \label{eq:norms}
\end{align}
For the middle estimate, $|A_P|\le3R^2$, $|B_P|\le12R^3$, and
$|C_P|\le3R^4$; subtracting two such quantities doubles each bound.
Multiplying the bounds for fifteen, sixty, and forty-five factors gives
\begin{equation}
 \|E_f\|_1\le4^{15}6^{165}R^{360}=:M(H).
 \label{eq:height-bound}
\end{equation}

\begin{theorem}[A search-free height separator]
\label{thm:height}
The positive integer
\begin{equation}
 t_\ast=4^{15}6^{165}(1+H)^{360}+2
 \label{eq:height-separator}
\end{equation}
separates both compressed families and both original descriptors on the
curve.
\end{theorem}
\begin{proof}
The nonzero polynomial $E_f$ has integer coefficients, so its leading
coefficient has absolute value at least one. Its Cauchy root bound is at
most $1+M(H)$ by \cref{eq:height-bound}. Hence $E_f(t_\ast)\ne0$.
Apply \cref{eq:exact-test} and the nonzero-multiplier identities.
\end{proof}

Although $t_\ast$ is numerically enormous, its bit length is $O(1+\log(1+H))$
with a large absolute constant. Its role is to show that even separator
selection itself can be replaced by an explicit arithmetic expression.
The first successful candidate in $1,\ldots,196$ will generally be much
smaller and avoid unnecessary coefficient growth.

\section{Exact checks, certificates, and their limits}
\label{sec:checks}

The attached computations are designed to test local algebra, combinatorial
counts, discriminant multiplicities, symmetry signs, and explicit examples.
They are not used as evidence that an unproved general statement should be
accepted. Each general theorem above has a written argument, including the
irreducibility and arbitrary-test-set sharpness results.

\subsection{Local symbolic identities}

The script \src{code/verify.py} uses SymPy to expand and check four universal
identities: the pair and triple descriptor expansions, the common-block
factorization, and the triple exchange factorization. These checks involve
independent root symbols, rather than approximate roots of sample
polynomials. They therefore test the precise signs and coefficients in
\cref{eq:pair-descriptor,eq:triple-descriptor,eq:shared,eq:exchange}.

\subsection{Integer arithmetic checks}

The same script enumerates the fifteen matchings and ten triple partitions
without a stored table. It checks the $45+60$ decomposition of matching
comparisons and the 45 disjoint-root-pair comparisons. For each root
transposition it checks that no edge of the block-disjoint matching graph
is reversed, matching the sign argument in \cref{lem:graph-sign}.

For four distinct integer root tuples and every $t=0,\ldots,16$, the
script checks both discriminant identities using exact integer products,
and verifies the compressed-to-descriptor identities and the exact
obstruction test. This yields 68 parameter cases, each testing both
resolvent discriminants. It also checks invariance of the three factors
under all fifteen root transpositions for those four tuples, giving
60 exact symmetry cases.

It exhausts all $\binom{11}{6}=462$ six-element subsets of
$\{-3,-2,\ldots,7\}$ and computes their first positive simultaneous
separator. The resulting distribution is
\begin{center}
\begin{tabular}{c*{8}{r}}
\toprule
First separator&1&2&3&4&5&6&7&8\\
\midrule
Number of tuples&93&83&89&125&48&15&8&1\\
\bottomrule
\end{tabular}
\end{center}
These finite data do not support a claim that eight works universally.
They provide reproducible examples, including the lower-bound witness in
\cref{prop:integer-lower}. Twenty scaled examples additionally check the
failure of fixed parameters described in \cref{prop:no-single}.

\subsection{A modular certificate for the sharp example}

The independent standard-library script
\src{code/sharpness_certificate.py} constructs $E_f$ modulo $1{,}000{,}003$
for \cref{eq:sharp-roots}. It uses only its 120 explicit factors, computes
an extended Euclidean identity with its derivative, and checks
\cref{eq:bezout}. The certificate stores every coefficient in ascending
order, so an independent implementation can verify the identity by
multiplication and addition without repeating the Euclidean algorithm.

Because reduction preserves the degree and gives gcd one, the
characteristic-zero polynomial is squarefree. Equivalently, the resultant
with its derivative is nonzero modulo the prime, hence is a nonzero
integer. This proves that the example reaches 195 distinct nonzero
complex failures for this parameterization. It does not imply that any of
those failures, let alone the first 195 positive integers, are all
integers or all real.

\subsection{Verification status}

Both scripts were executed successfully for this report. The package
records the exact results, the complete modular identity, and the scripts
needed to reproduce them. There was no Lean or Rocq execution for the new
results, no rebuild of the repository, and no benchmark against production
Galois software. The code supplied is an exact verifier and root-level
prototype, not a completed coefficient-only radical solver.

\section{A precise route back into ProveIt}
\label{sec:integration}

All repository statements here refer to snapshot \snapshot. The relevant
source directory is
\src{Algebra/PolynomialFormulas/Lean/PolynomialFormulas/}.
A subsequent repository version may have changed its interfaces or already
closed some of these obligations.

\subsection{Reusable existing boundaries}

The source already provides descriptor injectivity and evaluated-resolvent
separability equivalences in \src{SexticSeparatingInvariants.lean}; the
irreducible decision combines the chosen parameters with the Galois
criterion in \src{SexticIrreducibleDecision.lean}. The final all-roots
predicate and both computable and primitive-recursive criterion interfaces
are in \src{SexticRadicalComputability.lean}. Monicization is separated
cleanly into \src{SexticRadicalSemantics.lean}
\cite{repo-invariants,repo-search,repo-computability,repo-semantics}.

The reducible branch also contains a smaller search issue:
\src{linearTotalCode} uses \src{Nat.find}, even though a finite symmetric
integer-code bound is already present in \src{linearWitnessB}. A complete
primitive-recursive integration must replace or bound that selector as
well, not stop after fixing the pair and triple search
\cite{repo-reducible}.

\subsection{Suggested module sequence}

A useful first module, tentatively
\src{SexticCompressedDescriptors.lean}, would define $A,B,C,a,b,h,k$ and
prove \cref{eq:pair-descriptor,eq:triple-descriptor}. These are finite ring
identities, followed by short applications of the existing descriptor
injectivity theorem. Its main public result should identify the current
parameter $x(0)=2t^2,x(1)=t$ with the compressed values through the common
nonzero affine transformation.

A second module, \src{SexticCollisionFactors.lean}, would formalize the
shared-block/exchange identities and the $45/60/45$ finite partition
counts. It can obtain the $195$ bad-value bound directly from a product
over the finite factors, without first formalizing generic irreducibility
or the full discriminant formulas. Thus the shortest route to the bounded
selector is substantially smaller than the whole article.

A third module, \src{SexticBoundedSeparatingSearch.lean}, would prove that
one of $1,\ldots,196$ works and implement a fixed-length list scan. Its
return specification should establish simultaneous injectivity, followed
by the existing resolvent correctness bridge. A total fallback value can
be present in the definition, but its unreachable status on separable
inputs must be proved rather than assumed.

The computational module should then prove primitive recursiveness of
each fixed universal coefficient evaluation, of the squarefreeness test,
and of the bounded selector. One legitimate shortcut is to specialize at
each of the 196 constant parameters and unroll the finite search, provided
the coefficient evaluators themselves have primitive-recursive proofs.
Merely replacing \src{Nat.find} by a bounded loop around a function known
only to be computable would not establish \src{Primrec}.

Finally, replace the linear-factor selector by its existing bounded
interval, reuse the primitive-recursive quintic branch, and compose the
final criterion. The independent Rocq route can use the same mathematical
bound and compact collision products, but no claim of a completed port is
made here.

\subsection{Separate proof products from optional refinements}

The exact discriminant identities, generic irreducibility, and the
195-failure sharpness example deserve their own modules. They make the
mathematical story stronger, but they are not prerequisites for a correct
bounded sextic decision. Keeping them separate would preserve a small
trusted dependency chain for the computational endpoint.

A sensible audit should report four distinct endpoints: the universal
mathematical separator theorem; the correctness of the bounded returned
parameter; a primitive-recursive coefficient function; and the final
all-roots radical predicate. Exact sample evaluations are a fifth product,
not a substitute for any of the first four.

\section{A reusable general finite-separation principle}
\label{sec:general}

The sextic construction is special in its low degrees and factorization,
but the idea of converting polynomial descriptors to a finite hitting set
has an elementary general form. The following statement explains why an
existence proof by avoiding finitely many polynomial equalities often
admits an explicit bounded counterpart.

\begin{proposition}[Kronecker-curve separation]
\label{prop:kronecker}
Let $F_1(u,v),\ldots,F_m(u,v)$ be distinct polynomials over a
characteristic-zero field, each having $u$-degree at most $a$ and
$v$-degree at most $b$. Put
\[
 G_i(t)=F_i(t,t^{a+1}),\qquad D=a+b(a+1).
\]
Then the $G_i$ are distinct, have degree at most $D$, and among any
$D\binom m2+1$ distinct parameter values there is one at which all their
values are different.
\end{proposition}
\begin{proof}
A monomial $u^iv^j$ becomes $t^{i+(a+1)j}$. For $0\le i\le a$, the
exponent uniquely determines $i,j$ by division with remainder by $a+1$.
Thus the substitution is injective on the stated polynomial space and
preserves nonzero differences. The product of the $\binom m2$ differences
is nonzero of degree at most $D\binom m2$, and the root bound applies.
\end{proof}

This is a standard monomial-encoding argument, not a novelty claim. Its
contrast with the sextic result is instructive. A general curve must
preserve all descriptor coefficients. Here we first subtract the
partition-independent part, use the coefficient relations to compress
the family, and then recognize common collision factors. The result is a
quadratic curve, a low-degree obstruction, and an exact generic factor
count rather than a large undifferentiated degree estimate.

For other resolvent actions the promising order of work is therefore:
identify a complete descriptor; isolate its varying coefficients; find a
low-degree curve faithful to those coefficients; organize collision
pairs by group orbits; and only then use a root bound. The graph-product
sign test and the one-orbit product lemma can be reused independently of
the sextic setting.

\section{Further research questions and priorities}
\label{sec:questions}

The results leave several concrete problems. They fall into three groups:
sharp arithmetic separation, the geometry of the obstruction, and
certified implementation. The distinction between the solved arbitrary-set
bound and the unresolved integer-interval bound should be retained in
all follow-up work.

\begin{question}[The optimal initial integer interval]
Determine $N_{\mathrm{int}}$, currently bracketed by $8\le
N_{\mathrm{int}}\le196$ in this report. A proof that eight suffices would
require new control over simultaneous integer roots of the three factors,
not just an improvement in generic complex root counting. Conversely,
construct a distinct root tuple for which $1,\ldots,8$ all fail.
\end{question}
A practical first experiment is an exact search over larger integer and
rational root sets, recording which factor accounts for each failed
parameter. The combinatorial pattern of overlaps may be more informative
than the maximum alone. Any proposed universal improvement should be
proved independently of a finite sample.

\begin{question}[Short fixed test sets not forming an interval]
What is the smallest cardinality of a fixed set
$S\subset\ZZ_{>0}$ that separates every distinct complex root tuple?
The sharpness theorem for arbitrary input-dependent complex test sets
does not answer this question. Can a sparse integer set outperform an
initial interval substantially?
\end{question}

\begin{question}[Irreducible rational inputs]
Do the best uniform bounds decrease when $f$ is required to be
irreducible over $\QQ$? For a prescribed positive integer $t_0$, can one
construct an irreducible rational sextic with $E_f(t_0)=0$? The split
counterexamples in this article deliberately do not settle this.
\end{question}
A route through rational points on the three collision hypersurfaces
would have to preserve irreducibility of the sextic, rather than merely
produce a generic point in root coordinates.

\begin{question}[Alternative low-degree curves]
Can a different scalar compression or a different curve in descriptor
space lower the maximum number of exceptional parameters below 195?
The proof of sharpness applies to the particular $h,k$ parameterization,
not to every possible separating invariant.
\end{question}
The common-block factorization suggests looking for a curve that collapses
more repeated comparisons without making an entire family identical.
Any such proposal must still distinguish all partitions for every distinct
root tuple, including special symmetric configurations.

\begin{question}[Geometry of the three components]
Describe the intersections, singularities, and branching of the three
irreducible hypersurfaces $Q=0$, $U=0$, and $L=0$ over the separable-sextic
coefficient locus. Which specializations cause degree loss, repeated
parameter roots, or intersections between the components?
\end{question}
The generic irreducibility theorem reduces this from an unspecified
collision locus to three explicit orbit products. Their normalizations
and projection monodromy are natural next invariants.

\begin{question}[Galois groups of the obstruction factors]
Determine the Galois groups over $\QQ(e_1,\ldots,e_6)$ of the degree-30,
degree-120, and degree-45 parameter polynomials. Their factorizations into
quadratics and linears over the generic sextic splitting field provide
strong constraints, but do not by themselves determine the full groups.
\end{question}

\begin{question}[Compact coefficient circuits]
Construct certified straight-line programs for the two resolvents or for
$E_f(t)$ that avoid expanding large symmetric-polynomial tables. Compare
Newton identities, traces in finite algebras, and modular reconstruction,
with exact operation counts and coefficient-height estimates.
\end{question}
A useful first target is not the entire final solver but an independently
checkable evaluator for the obstruction at a supplied integer parameter.

\begin{question}[End-to-end formal primitive recursiveness]
Implement the module sequence in \cref{sec:integration}, close the bounded
linear-selector issue, and obtain a kernel-checked \src{Primrec} theorem
for the complete coefficient criterion. Port the same bounded search to
the independent Rocq development.
\end{question}
The decisive artifact would be a verified executable Boolean and its
all-complex-roots specification, not only a classical existence theorem
for a suitable natural number.

\begin{question}[Modular certificates and practical performance]
Can a small collection of finite-field evaluations certify good
parameters and rational-root decisions substantially faster than direct
integer expansion? Establish explicit reconstruction bounds and fallback
conditions so that modular success is a certificate rather than a
probabilistic guess.
\end{question}

\begin{question}[Positive characteristic and field size]
Determine the precise characteristics and finite-field sizes for which
analogues hold. The compressed formulas and many collision identities
survive algebraically, but the multiplier $2t^2$, the distinctness of
integer parameters, and generic irreducibility arguments must be
revisited. Extension-field parameters may be necessary.
\end{question}

\begin{question}[Higher-degree block systems]
For partitions into equal-sized blocks in higher degrees, classify
collision orbits and search for compressed descriptors with similarly
small invariant products. Which actions admit a graph-product
factorization and irreducible orbit factors of manageable degree?
\end{question}

\begin{question}[From a decision to explicit radical expressions]
Use a certified rational resolvent root to recover a fixed block system
and then construct branch-compatible radical expressions. A Boolean
solvability certificate is not yet an explicit all-roots formula; the
recovery and radical-tower interfaces require separate proofs.
\end{question}

\section{Conclusion}

The repository's unbounded separating search can be replaced by a
mathematically explicit finite mechanism. The central reason is not a
general computability theorem, but a small amount of partition-dependent
information: three coefficients for pair partitions, two for triple
partitions, and a quadratic curve that preserves them. The resulting
collision factors have enough structure to yield an exact discriminant
factorization and a generic irreducibility theorem.

This resolves the bounded-search obligation at the written-proof level
and identifies a precise next formalization target. It also settles a
sharpness question for arbitrary complex test sets while leaving the more
arithmetic positive-integer problem genuinely open. Keeping these scopes
separate is necessary both for mathematical accuracy and for a useful
research program.

\appendix
\section{Reproducibility and package contents}
\label{app:package}

The archive contains the self-contained source \src{article.tex}, its
compiled \src{article.pdf}, \src{README.md}, a build script, two Python
scripts, the exact verification report, and the modular B\'ezout
certificate. Bibliographic entries are embedded in the TeX source; no
separate bibliography tool is required. A checksum file records the
packaged artifacts.

Compile the document by running
\begin{verbatim}
./build.sh
\end{verbatim}
from the package directory. The script uses \src{pdflatex} three times and
writes the PDF in the same directory. The source uses commonly available
TeX Live packages, including \src{newtx}, \src{microtype}, and
\src{tcolorbox}.

Run the checks with
\begin{verbatim}
python code/verify.py
python code/sharpness_certificate.py
\end{verbatim}
The first script needs SymPy for its four local symbolic identities. The
second uses only the Python standard library and the combinatorial and
integer-arithmetic functions defined in \src{verify.py}; importing that
module does not import SymPy or run its test suite. The certificate
polynomials use ascending coefficient order. The default scripts use
assertions for test failures and should be run without Python's
optimization flag.

\section{Source and claim ledger}
\label{app:ledger}

The following ledger separates historical dependencies, inspected source
status, written results, and executed checks.

\begin{longtable}{@{}p{0.28\textwidth}p{0.66\textwidth}@{}}
\toprule
Claim or artifact & Basis and status\\
\midrule
\endhead
General radical-solvability decidability & Classical; Landau--Miller proves
polynomial-time decidability. Not a new result of this article.\\[5pt]
Sextic pair/triple group criterion & Classical; proved again in
\cref{sec:galois} to expose the exact role of separation.\\[5pt]
Complete block descriptors & Inherited from the inspected repository;
injectivity reproved in \cref{lem:complete}.\\[5pt]
Quadratic/linear compression & Written proof in \cref{sec:compression};
local identities checked symbolically.\\[5pt]
Three invariant factors and discriminants & Written proofs in
\cref{sec:obstruction,sec:discriminants}; exact integer tests supplied.\\[5pt]
Generic irreducibility and 196 sharpness & Written proofs in
\cref{sec:irreducibility,sec:bound}; explicit sharp example additionally
has a modular certificate.\\[5pt]
Initial interval optimum & Not determined; this report proves only
$8\le N_{\mathrm{int}}\le196$.\\[5pt]
Primitive-recursive coefficient criterion & Mathematical construction and
proof in \cref{sec:algorithm}; not an executed Lean development or a
packaged extracted solver.\\[5pt]
Executable checks & Both attached scripts executed successfully; exact
outputs in \src{checks/}. No floating-point evidence is used.\\[5pt]
Historical novelty & Specific contributions are relative to the inspected
repository and consulted sources; no exhaustive priority certification.\\
\bottomrule
\end{longtable}

\begin{thebibliography}{99}
\addcontentsline{toc}{section}{References}

\bibitem{landau-miller}
Susan Landau and Gary L. Miller,
\emph{Solvability by radicals is in polynomial time},
Journal of Computer and System Sciences \textbf{30}(2) (1985), 179--208.
DOI: \href{https://doi.org/10.1016/0022-0000(85)90013-3}{10.1016/0022-0000(85)90013-3}.
Author's bibliographic record:
\url{https://www.cs.cmu.edu/~glmiller/Publications/b2hd-LaMi85.html}.

\bibitem{boswell-glasser}
C. Boswell and M. L. Glasser,
\emph{Solvable sextic equations}, 2005 preprint,
arXiv:math-ph/0504001.
\url{https://arxiv.org/abs/math-ph/0504001}.
The two imprimitive subgroup types are described on page 3.

\bibitem{fieker-klueners}
Claus Fieker and J\"urgen Kl\"uners,
\emph{Computation of Galois groups of rational polynomials},
LMS Journal of Computation and Mathematics \textbf{17} (2014), 141--158.
DOI: \href{https://doi.org/10.1112/S1461157013000302}{10.1112/S1461157013000302}.
\url{https://arxiv.org/abs/1211.3588}.

\bibitem{repo}
Vladimir Reshetnikov, \emph{ProveIt: Polynomial formulas and the
Abel--Ruffini obstruction}, repository documentation, snapshot
\snapshot, inspected September 28, 2026.
\url{https://github.com/VladimirReshetnikov/ProveIt/blob/e2b1f016a94102663f12b970e6dd434229f90d01/Algebra/PolynomialFormulas/README.md}.

\bibitem{repo-decision}
Vladimir Reshetnikov, \emph{SexticRadicalDecision.lean}, same snapshot.
\url{https://github.com/VladimirReshetnikov/ProveIt/blob/e2b1f016a94102663f12b970e6dd434229f90d01/Algebra/PolynomialFormulas/Lean/PolynomialFormulas/SexticRadicalDecision.lean}.

\bibitem{repo-search}
Vladimir Reshetnikov, \emph{SexticIrreducibleDecision.lean}, same snapshot.
\url{https://github.com/VladimirReshetnikov/ProveIt/blob/e2b1f016a94102663f12b970e6dd434229f90d01/Algebra/PolynomialFormulas/Lean/PolynomialFormulas/SexticIrreducibleDecision.lean}.

\bibitem{repo-invariants}
Vladimir Reshetnikov, \emph{SexticSeparatingInvariants.lean}, same snapshot.
\url{https://github.com/VladimirReshetnikov/ProveIt/blob/e2b1f016a94102663f12b970e6dd434229f90d01/Algebra/PolynomialFormulas/Lean/PolynomialFormulas/SexticSeparatingInvariants.lean}.

\bibitem{repo-semantics}
Vladimir Reshetnikov, \emph{SexticRadicalSemantics.lean}, same snapshot.
\url{https://github.com/VladimirReshetnikov/ProveIt/blob/e2b1f016a94102663f12b970e6dd434229f90d01/Algebra/PolynomialFormulas/Lean/PolynomialFormulas/SexticRadicalSemantics.lean}.

\bibitem{repo-computability}
Vladimir Reshetnikov, \emph{SexticRadicalComputability.lean}, same snapshot.
\url{https://github.com/VladimirReshetnikov/ProveIt/blob/e2b1f016a94102663f12b970e6dd434229f90d01/Algebra/PolynomialFormulas/Lean/PolynomialFormulas/SexticRadicalComputability.lean}.

\bibitem{repo-reducible}
Vladimir Reshetnikov, \emph{SexticReducibleDecision.lean}, same snapshot.
\url{https://github.com/VladimirReshetnikov/ProveIt/blob/e2b1f016a94102663f12b970e6dd434229f90d01/Algebra/PolynomialFormulas/Lean/PolynomialFormulas/SexticReducibleDecision.lean}.

\end{thebibliography}
\end{document}
